%%
%% CGO 2027 submission: "Objects Without Allocation".
%% Loop-carried scalar replacement of persistent records in FOL.
%%
%% NOTE: The venue name below is set to CGO; the dates and location are
%% placeholders (TBD) to be filled from the real call for papers.
%%
\documentclass[sigplan,review,anonymous]{acmart}

\usepackage{listings}
\usepackage{booktabs} %% For formal tables
\usepackage{seqsplit}

\lstset{
  language=Lisp,
  basicstyle=\small\ttfamily,
  breaklines=true,
  frame=single,
  numbers=left,
  numberstyle=\tiny,
  keywordstyle=\color{blue},
  commentstyle=\color{green!40!black},
  stringstyle=\color{red},
  showstringspaces=false
}

%%
%% \BibTeX command to typeset BibTeX logo in the docs
\AtBeginDocument{%
  \providecommand\BibTeX{{%
    \normalfont B\kern-0.5em{\scshape i\kern-0.25em b}\kern-0.8em\TeX}}}

%% Rights management information.
\setcopyright{acmlicensed}
\copyrightyear{2027}
\acmYear{2027}
\acmDOI{XXXXXXX.XXXXXXX}

%% These commands are for a PROCEEDINGS abstract or paper. (Dates/location TBD.)
\acmConference[CGO '27]{IEEE/ACM International Symposium on Code Generation and Optimization}{February 2027}{TBD}
\acmBooktitle{Proceedings of the IEEE/ACM International Symposium on Code Generation and Optimization (CGO '27)}
\acmPrice{15.00}
\acmISBN{978-1-4503-XXXX-X/27/06}

%%
%% The "title" command has an optional parameter for a short title.
\title[Objects Without Allocation]{Objects Without Allocation: Loop-Carried Scalar Replacement of Persistent Records in a Dynamic Language}

\author{Anonymous Author(s)}
\affiliation{%
  \institution{Anonymous Institution}
  \city{Anonymous City}
  \country{Anonymous Country}
}
\email{anonymous@example.com}

\renewcommand{\shortauthors}{Anonymous Author(s)}

\begin{abstract}
Persistent-first languages give programmers immutable objects, but a record
rebuilt every iteration of a loop turns into a stream of short-lived heap
allocations: on our system a loop that rebuilds a persistent record runs
90--290$\times$ slower than the same loop over a native mutable struct,
dominated by allocation and garbage collection. Classic
\emph{scalar replacement} removes such allocations by holding an object's fields
in registers, but the textbook formulation only fires when an object is created
and dies inside one lexical scope---the least profitable case---and it assumes a
closed world in which a class's shape never changes. We present \emph{aggregate
scalar replacement}, an emit-time transformation for FOL (a persistent-first
Lisp) that unboxes a persistent-record \emph{loop accumulator}: the accumulator
becomes one scalar loop variable per field, field reads become variable reads,
the reconstruction that feeds the back-edge becomes per-field value expressions,
and a single object is re-boxed only at loop exit. Two ingredients make this
practical in a dynamic language: (1) an \emph{inlining} step that exposes the
constructor returned by an object-producing callee at the recur position, giving
interprocedural reach without a whole-program analysis; and (2) a
\emph{world-guarded} dual path that keeps the optimization sound when a class is
redefined at run time. A single classify-and-rewrite walk is aligned by
construction---any use it does not recognize aborts the rewrite and yields the
original loop---so the transformation is never unsound, only sometimes
inapplicable. Across four allocation-bound benchmarks the technique eliminates
per-iteration allocation entirely for an end-to-end 5.6$\times$--7.2$\times$
speedup, recovering the allocation component of that 90--290$\times$ gap, with
byte-identical results and no cost when disabled; the residual distance to
native, which we measure and attribute to generic scalar arithmetic rather than
allocation, is orthogonal to this work.
\end{abstract}

\begin{CCSXML}
<ccs2012>
   <concept>
       <concept_id>10011007.10010940.10010941.10010942.10010943</concept_id>
       <concept_desc>Software and its engineering~Compilers</concept_desc>
       <concept_significance>500</concept_significance>
       </concept>
   <concept>
       <concept_id>10011007.10011006.10011041.10011045</concept_id>
       <concept_desc>Software and its engineering~Dynamic compilers</concept_desc>
       <concept_significance>500</concept_significance>
       </concept>
   <concept>
       <concept_id>10003752.10003753.10003757</concept_id>
       <concept_desc>Theory of computation~Program analysis</concept_desc>
       <concept_significance>500</concept_significance>
       </concept>
   <concept>
       <concept_id>10002944.10011122.10011123</concept_id>
       <concept_desc>General and reference~Performance</concept_desc>
       <concept_significance>300</concept_significance>
       </concept>
 </ccs2012>
\end{CCSXML}

\ccsdesc[500]{Software and its engineering~Compilers}
\ccsdesc[500]{Software and its engineering~Dynamic compilers}
\ccsdesc[500]{Theory of computation~Program analysis}

\keywords{scalar replacement, escape analysis, persistent data structures,
inlining, dynamic languages, lisp, compiler optimization, speculative
optimization, deoptimization}

\begin{document}
\maketitle

\section{Introduction}

Languages built on persistent data structures let programmers treat records as
immutable values: an ``update'' produces a new version that structurally shares
with the old one. This is a pleasant programming model, but it has a cost. When
a small record is updated inside a loop, every iteration allocates a fresh
object on the heap, and the previous version becomes garbage immediately. In
FOL---a Lisp-1 dialect with CLOS-style persistent objects, transpiled to Common
Lisp---a loop that rebuilds a persistent record runs \textbf{90--290$\times$}
slower than the same loop written over a native mutable \texttt{defstruct}
(\S\ref{sec:eval}), dominated by allocation and garbage-collection pressure.

The canonical remedy in a static, closed-world setting is \emph{scalar
replacement of aggregates} (SRA): if a compiler can prove an object does not
escape, it need not allocate the object at all; the object's fields live in
local variables, and field accesses become variable accesses. Modern
just-in-time compilers realize a powerful form of this via \emph{partial escape
analysis} \cite{stadler2014partial}, building on a long line of escape analyses
for object-oriented languages \cite{choi1999escape,blanchet2003escape}.

Two obstacles keep the textbook formulation from helping the programs that hurt
most in a language like FOL. First, the \emph{scope} problem: the simplest SRA
only fires when an object is created and dies inside a single lexical
scope---a \texttt{let} whose body consumes the object and returns something
else. That is the least interesting case; a good backend often stack-allocates
it anyway. The objects that actually drive allocation are \emph{loop
accumulators}: a record that is read, rebuilt, and threaded through a
\texttt{recur} back-edge, iteration after iteration, frequently by calling a
small helper that returns a fresh record. Second, the \emph{open-world}
problem: FOL classes can be redefined at run time. An optimization that bakes in
a class's field layout must remain correct even if that class's shape changes
\emph{while the optimized code is running}---something a closed-world SRA never
has to consider.

This paper presents \emph{aggregate scalar replacement} for FOL: an emit-time
transformation that unboxes a persistent-record loop accumulator into one scalar
loop variable per field. Figure~\ref{fig:overview} shows the target pattern and
the result. Our contributions are:

\begin{enumerate}
  \item \textbf{Loop-carried record unboxing.} A transformation that splits a
  record accumulator into per-field scalar loop variables: \texttt{(get p :x)}
  becomes a variable read, the constructor that feeds \texttt{recur} becomes its
  field value-expressions (the loop's arity grows to match), and a single
  object is re-materialized only if the bare accumulator is returned at the loop
  exit (\S\ref{sec:loop}).
  \item \textbf{Interprocedural reach by inlining.} An object-returning callee
  at the recur position---the common shape \texttt{(recur (update p) \ldots)}---is
  inlined so that its tail constructor becomes the reconstruction expression,
  giving interprocedural unboxing without a whole-program ownership analysis
  (\S\ref{sec:inline}).
  \item \textbf{Soundness under live redefinition.} The converted loop is
  emitted as a \emph{world-guarded} dual path keyed on the record's class name;
  redefining the class invalidates the region and the next call transparently
  takes the original, object-allocating loop (\S\ref{sec:world}).
  \item \textbf{Alignment by construction.} A single walk both classifies and
  rewrites; any accumulator use it does not recognize aborts the rewrite and
  yields the original loop unchanged. The transformation is therefore never
  wrong, only sometimes inapplicable---the safe default for a compiler pass
  enabled on untyped code (\S\ref{sec:align}).
  \item \textbf{Evaluation.} On four allocation-bound benchmarks---spanning
  two- and three-field records, additive and multiplicative update shapes, and a
  coupled two-accumulator loop---the technique reduces per-iteration allocation
  to zero for an end-to-end 5.6$\times$--7.2$\times$ speedup at 5M iterations,
  produces
  byte-identical results, and adds no cost when disabled or inapplicable. We
  additionally measure a native mutable-\texttt{defstruct} implementation of
  each benchmark as an upper bound, and report honestly how much of the gap ASR
  closes and what the residual is (\S\ref{sec:eval}).
\end{enumerate}

\begin{figure}
\begin{lstlisting}
;; --- Source ---
(defclass <point> [] [[x] [y]])

(defn update-point [p]
  (bind [dx 0.1 dy 0.2]
    (make-<point> :x (+ (get p :x) dx)
                  :y (+ (get p :y) dy))))

(defn run-simulation [n]
  (loop [p (make-<point> :x 0 :y 0) i 0]
    (if (< i n)
      (recur (update-point p) (inc i))
      p)))

;; --- After aggregate scalar replacement ---
(loop [p_x 0 p_y 0 i 0]
  (if (< i n)
    (bind [dx 0.1 dy 0.2]
      (recur (+ p_x dx) (+ p_y dy) (inc i)))
    (make-<point> :x p_x :y p_y)))
\end{lstlisting}
\caption{The target pattern. \texttt{update-point} is inlined at the recur
position, the accumulator \texttt{p} is split into scalar loop variables
\texttt{p\_x}/\texttt{p\_y}, and the object is re-boxed only at the exit. In
the emitted code this fast path sits behind a world guard on \texttt{<point>}
(Figure~\ref{fig:emitted}).}
\label{fig:overview}
\end{figure}

Our evaluation is deliberately focused: the technique targets one high-value
pattern---the loop-carried record accumulator---and we measure it on four
microbenchmarks that instantiate that pattern with different field counts, update
shapes, and accumulator multiplicities. We make no claim here about how
frequently the pattern arises in
production code; a corpus study of its incidence across real FOL programs is
future work (\S\ref{sec:eval}). What we do report is an honest picture of when
the transformation fires and what it buys where it does.

\section{Background}

\textbf{FOL in one column.} FOL is a Lisp-1 dialect transpiled to Common Lisp
(SBCL). Persistent records are declared with \texttt{defclass}; a class
\texttt{<t>} induces a constructor \texttt{make-<t>} and keyword field
accessors read with \texttt{(get obj :field)}. A global type registry records,
for each class, its field set. \texttt{(loop [bindings] \ldots (recur
new-vals))} provides tail-recursive iteration; the compiler lowers it to a
\texttt{block}/\texttt{tagbody} with a \texttt{psetq} that updates all loop
variables simultaneously before jumping to the head. Forms are transpiled one
at a time, in source order.

\textbf{Why the accumulator is the hard case.} In
Figure~\ref{fig:overview}, the allocation lives in two places, neither of which
a scope-local SRA can touch. In \texttt{update-point}, the \texttt{make-<point>}
is the function's \emph{return value}---it escapes to the caller. In
\texttt{run-simulation}, \texttt{p} is a \emph{loop parameter}, not a
\texttt{let}-binding, and it is passed whole into \texttt{update-point}. To
remove the allocation one must reason across the call boundary and across the
loop back-edge simultaneously. That is precisely what aggregate scalar
replacement does.

\textbf{Relationship to transient conversion.} A companion line of work makes
persistent \emph{collection} accumulators (dicts, vectors, sets) update in place
by converting them to Clojure-style transients, guarded for soundness under
redefinition. Aggregate scalar replacement is the analogue for user-defined
\emph{records}: rather than mutate a boxed object in place, it eliminates the box
entirely. The two share the world-guard machinery described in
\S\ref{sec:world} and compose in the same \texttt{emit-loop} entry point. That
machinery is shared infrastructure, not a claim of this paper: our contribution
is the record transformation of \S\ref{sec:algo}, and we describe the guard in
full only because it is what makes that transformation sound in an open world.
The transient conversion is separate, concurrent work, mentioned here solely for
context.

\section{Aggregate Scalar Replacement}
\label{sec:algo}

The transformation runs at loop-emission time. Given a \texttt{loop} node, it
searches the bindings for a parameter whose initializer is a record constructor
\texttt{(make-<t> \ldots)} of a class known to the type registry, and attempts
to unbox it. The attempt either succeeds and returns a rewritten loop, or bails
and leaves the loop untouched. We describe the transform for a single
accumulator; a loop with several is handled by applying it to a fixpoint
(\S\ref{sec:discussion}). We cover candidate selection, the two core rewrites,
the alignment discipline that keeps them sound, and the world guard.

\subsection{Candidates and fields}
\label{sec:cand}

A loop parameter \texttt{p} is a \emph{candidate} when (i) its initializer is a
constructor call \texttt{(make-<t> :k1 v1 \ldots)}, (ii) \texttt{<t>} is in the
type registry with a known field set $F = \{f_1,\ldots,f_m\}$, and (iii) the
constructor supplies exactly the fields in $F$. For each field $f_j$ we mint a
fresh scalar loop variable named \texttt{p\_}$f_j$ (e.g.\ \texttt{p\_x}). The
loop binding for \texttt{p} is replaced by $m$ bindings \texttt{p\_}$f_j$
initialized to the corresponding constructor argument $v_j$. All keys are
canonicalized so that the constructor initarg, the \texttt{get} key, and the
registry's storage key coincide.

\subsection{Loop-carried unboxing}
\label{sec:loop}

The body is rewritten by a single recursive walk parameterized by whether the
current node is in tail position and by a set of \emph{alias} names that denote
the accumulator (initially $\{\texttt{p}\}$; \S\ref{sec:inline} explains why it
is a set). The walk maps each accumulator use to one of a small number of
recognized shapes (Figure~\ref{fig:algo} gives it in full):

\begin{itemize}
  \item \textbf{Field read.} \texttt{(get p :}$f_j$\texttt{)} rewrites to a
  reference to \texttt{p\_}$f_j$. A read of a key not in $F$ aborts.
  \item \textbf{Reconstruction at the back-edge.} At the recur position, the
  argument in the accumulator's slot must \emph{reconstruct} the record. The
  walk expands it into the $m$ per-field value expressions and splices them into
  the recur argument list in field order, growing the loop's arity from
  $k$ to $k+m-1$. A bare accumulator here (an unchanged pass-through) expands to
  the current field variables.
  \item \textbf{Re-box at exit.} A bare accumulator reference in tail position
  (the loop's result) is re-materialized as \texttt{(make-<t> :}$f_1$
  \texttt{p\_}$f_1$ \ldots\texttt{)}---a single allocation, executed at most once
  per loop activation.
  \item \textbf{Anything else aborts.} A bare accumulator passed to an
  arbitrary call, stored in a literal, captured by a closure, or reaching a node
  type the walk does not handle causes the whole attempt to fail. Because the
  walk builds the rewritten body functionally and signals failure by a
  non-local exit, a failed attempt has no effect: the original loop is emitted.
\end{itemize}

\begin{figure}
\begin{lstlisting}[language={},basicstyle=\footnotesize\ttfamily,numbers=none,frame=single,columns=fullflexible,keepspaces=true,keywordstyle={}]
Rewrite(n, tail, A):        # A = accumulator alias names; may raise Fail
  match n:
    (get p k), p in A    -> if k in F: return VAR(p_k) else Fail
    (recur ... a ...)    -> replace slot arg a by Reconstruct(a, A),  # k -> k+m-1
                            other args via Rewrite(., false, A)
    p, p in A and tail   -> return (make-<T> :f1 VAR(p_f1)...:fm VAR(p_fm)) # re-box
    p, p in A            -> Fail                    # bare accumulator escapes
    (bind bs e)|(do e)   -> peel; rewrite inits (false) and body (tail); re-wrap
    _                    -> recurse; any surviving p in A -> Fail

Reconstruct(a, A):          # returns m per-field expressions, in field order
  a := Peel(a)                              # strip single-form bind/do wrappers
  match a:
    (make-<T> :f1 v1...:fm vm), {f_j} = F
                         -> return [ Rewrite(v_j, false, A) | j = 1..m ]
    (g x1...xn), g in Inlinable, every x_i a symbol or literal
                         -> (ps, body) := Inlinable[g]; q := param of g bound to acc
                            return Reconstruct(body, A U {q}) # other ps -> wrapper
    p, p in A            -> return [ VAR(p_f1)...VAR(p_fm) ]  # unchanged pass-through
    _                    -> Fail
\end{lstlisting}
\caption{The classify-and-rewrite walk (\S\ref{sec:loop},~\S\ref{sec:inline}).
$F=\{f_1,\dots,f_m\}$ is the accumulator class's field set, \texttt{VAR}($p\_f$)
the scalar loop variable for field $f$, and \texttt{Inlinable} the table mapping
a single-clause constructor-returning function to its parameters and body
(\S\ref{sec:inline}). Classification and rewriting are one traversal: any
accumulator use not matched here raises \texttt{Fail}, whereupon the original
loop is emitted unchanged.}
\label{fig:algo}
\end{figure}

A reconstruction may be wrapped in single-form \texttt{bind}/\texttt{do} layers
(as in Figure~\ref{fig:overview}, where the \texttt{make-<point>} sits under
\texttt{(bind [dx 0.1 dy 0.2] \ldots)}). The walk peels these layers, rewrites
their initializers (which may themselves read accumulator fields), and re-wraps
the rewritten \texttt{recur} in the same layers, so that names like \texttt{dx}
remain in scope for the field expressions that use them.

Because the emitted \texttt{recur} lowers to a \texttt{psetq} over the loop's
(now $m$-way expanded) binding list, and both the binding list and every
\texttt{recur} are rewritten together, the loop stays internally consistent by
construction; no change to the loop/recur code generator is required.

\subsection{Interprocedural reach by inlining}
\label{sec:inline}

The accumulator rarely reconstructs itself with a literal constructor; more
often it is rebuilt by a helper, as \texttt{(recur (update-point p) \ldots)}. To
expose the helper's constructor, we inline it. During transpilation, a
lightweight \texttt{compile-form} hook records every single-clause function whose
body reduces (through \texttt{bind}/\texttt{do} peeling) to a record
constructor into an \emph{inlinable} table, keyed by name.

When the reconstruction at a recur position is a call to an inlinable function
whose accumulator-carrying argument is the accumulator itself, the walk inlines
the callee's body in place. The parameter that received the accumulator is added
to the alias set, so that reads of \emph{that} parameter inside the inlined body
(e.g.\ \texttt{(get p :x)} in \texttt{update-point}, whose parameter is also
named \texttt{p}) are recognized as accumulator reads and rewritten to field
variables. Other parameters are bound by a wrapper layer around the rewritten
\texttt{recur}. Inlining is applied only when the callee's arguments are
symbols or literals, which keeps the substitution capture-free and avoids
duplicating side effects. This is deliberately a one-level, syntactic inliner:
it is enough to reach the dominant \texttt{(recur (helper acc) \ldots)} shape,
and it composes with the peeling of \S\ref{sec:loop} to turn Koka-style
``functional but in place'' record threading into flat scalar arithmetic.

Unlike ownership- or effect-based approaches to in-place update
\cite{reinking2021perceus,lorenzen2023functional}, this reach requires no type
information: the inlinable table is populated purely from syntactic shape, and
the alias-set discipline keeps the rewrite local and checkable.

\subsection{Alignment by construction}
\label{sec:align}

The classifier and the rewriter are the \emph{same walk}. There is no separate
``is this legal?'' analysis that could drift out of sync with the
transformation: the walk rewrites what it recognizes and aborts on everything
else. This has two consequences. First, soundness of the field-level rewrite is
local and obvious: every surviving accumulator reference is either a field read
(replaced by the corresponding scalar) or a tail re-box (replaced by an
equivalent constructor), and no bare accumulator escapes---if one did, the walk
would have aborted. Second, the pass is conservative by default: novel or
complex accumulator uses simply prevent optimization rather than risk a wrong
rewrite. For a pass enabled on latently-typed code with open-world dispatch,
``inapplicable'' is a far cheaper failure mode than ``incorrect.''

We can state the guarantee precisely. Say \texttt{Rewrite} (Figure~\ref{fig:algo})
\emph{succeeds} on a loop body if it returns without raising \texttt{Fail}.

\smallskip
\noindent\textbf{Invariant (alignment).} \emph{If \texttt{Rewrite} succeeds, then
in the rewritten body every reference to an accumulator-alias name has been
replaced by one of exactly three forms: (i) the scalar variable
\texttt{VAR}($p\_f$) of a field $f\in F$ (a field read); (ii) the per-field value
expressions spliced into a \texttt{recur} (a reconstruction); or (iii) the single
tail re-box \texttt{(make-<t> \ldots)}. No bare accumulator reference survives.}

\smallskip
\noindent The invariant holds by construction: those three are the only cases in
which \texttt{Rewrite} returns a term mentioning the accumulator, and every other
occurrence---a field read of an absent key, a bare accumulator in non-tail
position, an unrecognized reconstruction---is a \texttt{Fail} case. Two
corollaries follow. \emph{Allocation:} the only \texttt{make-<t>} reachable in the
rewritten loop is the exit re-box of case (iii), which post-dominates the loop and
runs at most once per activation, so the back-edge allocates nothing.
\emph{Correctness:} because the walk is total---it either succeeds with the
above shape or emits the original loop verbatim---the pass never yields a rewrite
that is not one of these three equivalences, which is why it is safe to enable on
untyped code without a separate legality analysis.

\subsection{Soundness under live redefinition}
\label{sec:world}

Aggregate scalar replacement bakes a class's field layout into the fast path: it
assumes \texttt{<point>} has exactly fields \texttt{x} and \texttt{y}. In FOL a
class can be redefined at run time---\texttt{(defclass <point> [] [[x] [y]
[z]])}---which would invalidate that assumption. We keep the optimization sound
with the same \emph{invalidation ladder} the transient work uses---a static,
ahead-of-time analogue of the assumption-and-deoptimization mechanism JIT
compilers use for speculative optimization (\S\ref{sec:related}).

The converted loop is emitted as a guarded dual path
(Figure~\ref{fig:emitted}). At load time, a validity cell is registered and
indexed under the record's class name. \texttt{emit-defclass} is instrumented so
that (re)defining a class \emph{notifies} the world, flipping the cells of
exactly the regions that assumed that class name. The guard---one pointer read
of the cell's \texttt{car}, checked once on loop entry---selects the unboxed fast
path while the assumption holds and the original object-allocating loop after
the class is redefined. In-flight activations complete on the path they entered,
which is safe because the fast path performs only scalar arithmetic and the
final re-box against the field layout it was compiled for; we make the
concurrent execution model of this guard precise in \S\ref{sec:discussion}. When
the compiler is told the world is sealed (whole-program, cold-loaded builds), the
guard is omitted entirely and only the fast path is emitted.

\begin{figure}
\begin{lstlisting}
(IF (CAR (LOAD-TIME-VALUE
           (REGISTER-REGION '("<POINT>")) T))
    ;; fast path: accumulator unboxed, callee inlined
    (BLOCK LOOP-BLOCK-2
      (LET ((P_X 0) (P_Y 0) (I 0))
        (TAGBODY LOOP-2
          ... (IF (< I ITERATIONS)
                  (LET ((DX 0.1)) (LET ((DY 0.2))
                    (PSETQ P_X (+ P_X DX)
                           P_Y (+ P_Y DY)
                           I   (INC I))
                    (GO LOOP-2)))
                  (MAKE-<POINT> :X P_X :Y P_Y)) ...)))
    ;; fallback: original allocating loop
    (BLOCK LOOP-BLOCK-3
      (LET ((P (MAKE-<POINT> :X 0 :Y 0)) (I 0))
        ... (PSETQ P (UPDATE-POINT P) I (INC I)) ...)))
\end{lstlisting}
\caption{Emitted Common Lisp for \texttt{run-simulation}. The world guard reads
one cell registered under \texttt{"<POINT>"}; redefining \texttt{<point>} flips
the cell and the next call takes the fallback. The fast path contains no
constructor call on the back-edge---only \texttt{psetq} of scalar arithmetic.}
\label{fig:emitted}
\end{figure}

\section{Implementation}
\label{sec:impl}

The pass is implemented inside FOL's compiler in about 300 lines of Common Lisp.
It lives at the loop-emission boundary rather than as a standalone AST pass, for
a concrete reason: the transformation needs the class registry and the record
constructors that \texttt{defclass} populates during emission, and it reuses the
existing loop/recur code generator by handing it an already-rewritten loop node.
Its entry point is consulted by \texttt{emit-loop} before the transient path;
when it declines, emission proceeds exactly as before, so a program that never
triggers the pattern produces identical output.

The inlinable-function table is populated by the same \texttt{compile-form} hook
that drives interprocedural summary inference, and because forms are transpiled
in source order, a helper is registered before the loop that calls it. The pass
is gated behind a single flag; with the flag off, \texttt{emit-loop} is
byte-for-byte unchanged.

\section{Evaluation}
\label{sec:eval}

We evaluate four questions. \textbf{RQ1:} does the transformation fire on the
target pattern and eliminate allocation? \textbf{RQ2:} what speedup does it
deliver, and is the result identical to the unoptimized program? \textbf{RQ3:}
what does it cost when disabled or inapplicable? \textbf{RQ4:} how close does the
optimized code come to a native mutable-struct implementation, and what accounts
for the difference? All measurements were taken on an AMD Ryzen 5 7430U (used
single-threaded) with 32\,GB of RAM, running SBCL 2.6.0---generational GC,
1\,GB dynamic space, 51\,MB nursery---on Windows 11. For each benchmark and
configuration we run one warm-up call followed by five timed calls, forcing a
full collection before each, and report the mean $\pm$ sample standard
deviation. Wall time uses the real-time clock; per-call allocation,
garbage-collection counts, and time spent in GC are read from the SBCL runtime
(the bytes-consed counter, an after-GC hook, and \texttt{*gc-run-time*});
speedup is computed per paired trial and then averaged. As an upper bound we
also implement each benchmark natively---an idiomatic mutable \texttt{defstruct}
updated in place, with declared \texttt{single-float} slots matching FOL's float
type, compiled at \texttt{(speed 3) (safety 0)}. Because a speedup over the
baseline depends on how often the collector runs, it carries run-to-run
variation of roughly $\pm0.5\times$; the allocation figures (deterministic) and
the native ceiling (nearly noise-free) are the more robust numbers. The four
benchmarks all instantiate the loop-carried record accumulator but differ in
shape: a \emph{particle} simulation (two independent additive field updates, the
pattern of Figure~\ref{fig:overview}), a 2D \emph{rotation} orbit (two fields
updated by a coupled, multiplicative map), a three-field \emph{ballistic}
integrator (three fields, where the position update reads the freshly computed
velocity), and a coupled \emph{two-body} relaxation that carries \emph{two}
\texttt{<vec2>} record accumulators, each step reading both---the
multi-accumulator case.

\textbf{RQ1: Firing and allocation.} The pass fires on all four benchmarks: each
emitted fast path (Figure~\ref{fig:emitted}) splits every record accumulator into
one scalar loop variable per field, inlines any update helper, and contains
\emph{no} constructor call on the back-edge. In the two-body benchmark the
fixpoint (\S\ref{sec:discussion}) unboxes \emph{both} accumulators---four scalar
variables---and the parallel \texttt{recur}/\texttt{psetq} keeps the coupled
update correct. The effect is visible directly in the runtime counters: each
baseline conses 1.1--2.2\,GB per call and triggers 20--42 collections costing
25--106\,ms, whereas the scalar-replaced version conses \emph{nothing} and
triggers no workload collection at all. Per-iteration allocation drops from one
record (224 or 256\,B), or two records (448\,B) in the two-body case, to zero;
the only allocation that survives is the single re-box at loop exit, executed at
most once per activation. The native versions, mutating their structs in place,
likewise cons nothing per iteration.

\textbf{RQ2: Speedup and correctness.} Table~\ref{tab:speedup} reports the four
benchmarks at five million iterations. Eliminating allocation removes the
garbage-collection pressure that dominates each baseline, yielding end-to-end
speedups of 5.6$\times$--7.2$\times$. Crucially, every optimized program returns
results bit-identical to its unoptimized counterpart across all five trials; a
dedicated test suite checks optimized-equals-unoptimized on several
record-accumulator shapes (field read at exit, bare re-box at exit, inlined and
non-inlined reconstructions, and multiple coupled accumulators), and the full
compiler test suite of 3{,}299 checks passes with the pass enabled.

\begin{table}
  \caption{Four allocation-bound benchmarks at 5M iterations: wall time (mean
  $\pm$ sample stddev over five timed trials, after one warm-up) for the
  persistent baseline, aggregate scalar replacement (ASR), and a native
  mutable-\texttt{defstruct} upper bound; AMD Ryzen 5 7430U, SBCL 2.6.0. ASR
  drives per-iteration allocation (224--448\,B) to zero and is bit-identical to
  the baseline; the ``$\times$'' columns are speedups over the baseline.
  Two-body carries \emph{two} coupled record accumulators, both unboxed.}
  \label{tab:speedup}
  \small
  \begin{tabular}{lrrrrr}
    \toprule
    Benchmark & Base & ASR & Native & ASR & Native \\
              & (ms) & (ms) & (ms) & $\times$ & $\times$ \\
    \midrule
    Particle  & $1908\pm46$ & $294\pm1$  & 15 & 6.5 & 128 \\
    Rotation  & $2745\pm76$ & $492\pm11$ & 30 & 5.6 & \phantom{0}93 \\
    Ballistic & $2514\pm45$ & $371\pm3$  & 15 & 6.8 & 168 \\
    Two-body  & $6951\pm41$ & $964\pm5$  & 24 & 7.2 & 294 \\
    \bottomrule
  \end{tabular}
\end{table}

\textbf{RQ4: Distance to native.} A hand-written native version---a mutable
\texttt{defstruct} updated in place---is 93$\times$--294$\times$ faster than the
persistent baseline (Table~\ref{tab:speedup}), which is what makes the pattern
worth attacking in the first place. Aggregate scalar replacement recovers the
\emph{allocation} component of that gap: it removes 100\% of per-iteration
consing, worth the 5.6$\times$--7.2$\times$ of RQ2. It does \emph{not} close the
gap entirely---the scalar-replaced loop remains 16$\times$--41$\times$ slower
than native, reaching 2--6\% of native speed---and we report that rather than
eliding it. The gap is widest on the two-body benchmark, whose heavier per-step
arithmetic makes the generic-operator overhead (below) proportionally larger.
The residual is not allocation, of which none remains, but
arithmetic: the unboxed loop still calls FOL's \emph{generic} operators
(\texttt{+}, \texttt{*}, \texttt{-}) on its scalar variables, whereas the native
loop issues declared single-float machine instructions. Closing it is a matter
of scalar type specialization, which is orthogonal to allocation and in fact
complementary to this transformation---scalar replacement is what exposes the
per-field variables that a type specializer would then annotate. As an
allocation optimization, ASR is complete on these benchmarks; the remaining
distance is a separate axis of work.

\textbf{RQ3: Cost when disabled or inapplicable.} With the flag off, loop
emission is unchanged and the output is byte-identical to the unmodified
compiler; there is no residual cost. When the flag is on but the pattern does
not apply---because the accumulator escapes, the reconstruction is not a
recognized shape, or the class is unknown---the walk aborts and the original
loop is emitted, again with no change to the generated code. We verified the
abort path on an accumulator that flows into an unknown call: the loop is left
untouched and no scalar variables are introduced.

\textbf{Correctness under redefinition.} We exercise the world guard directly:
after a converted loop registers its region, redefining the record class flips
the region's validity cell, and the next call falls back to the original
allocating loop and still produces the correct result against the redefined
class.

\section{Discussion and Limitations}
\label{sec:discussion}

\textbf{Scope of the current transformation.} The pass unboxes \emph{every}
qualifying record accumulator in a loop, applying the single-accumulator rewrite
to a fixpoint: it unboxes one, re-scans the rewritten loop for another, and
repeats until none remain. Because \texttt{recur} lowers to a parallel
\texttt{psetq}, coupled accumulators---where one's reconstruction reads another's
fields---remain correct, each cross-read resolving to the other's scalar
variables as it is unboxed in turn; an accumulator that fails to qualify is left
boxed, so replacement is per-accumulator rather than all-or-nothing (the
two-body benchmark unboxes two coupled records this way). Non-record loop
variables (counters, indices) pass through unaffected. Genuine limitations
remain: each reconstruction must peel---through single-form
\texttt{bind}/\texttt{do} wrappers and at most one level of inlinable call---to a
constructor that supplies \emph{every} field; \texttt{if}-branched
reconstructions (\texttt{(if c (make-<t> \ldots) (make-<t> \ldots))}) are not yet
recognized and cause a conservative abort, as do inlinable callees whose
accumulator argument is a non-trivial expression or that return more than one
record. These are engineering extensions rather than fundamental limits: each
corresponds to a shape the walk could learn to peel, and each preserves the
safe-by-default discipline because an unrecognized shape aborts.

\textbf{Relationship to partial escape analysis.} Partial escape analysis with
scalar replacement \cite{stadler2014partial} is strictly more general: it
handles arbitrary control flow and materializes objects lazily at the points
where they escape. Our transformation trades that generality for a simple,
auditable rewrite specialized to the loop-accumulator pattern, and it adds
something a closed-world JIT does not need---soundness under class
redefinition. A promising direction is to generalize the single classify-rewrite
walk toward a materialization-on-escape discipline while retaining the world
guard.

\textbf{Freshness, not linearity.} Where Perceus \cite{reinking2021perceus} and
uniqueness/borrowing systems \cite{lorenzen2023functional} establish in-place
reuse through reference counting or types, we establish it structurally: the
accumulator is freshly constructed at loop entry and every use is either a field
read or a tail reconstruction, so no live alias survives a rewrite. This buys
in-place behavior---indeed, no ``place'' at all---without a type system, at the
cost of applying only where the syntactic discipline holds.

\textbf{Execution model and concurrency.} The world registry is shared mutable
state, so we state its contract precisely. All \emph{writers}---region
registration and the invalidation triggered by a (re)definition---are serialized
by a single lock. The \emph{reader} is the emitted guard: a deliberately
unsynchronized single-word read of the cell's validity flag, kept lock-free
because it sits on the hot path. Soundness of that unsynchronized read rests on
two properties of the flag. It is \emph{monotonic}: the transition is one-way,
$\text{valid}\to\text{invalid}$, and an invalidated cell is unlinked from the
dependency map and never revived; and its write is a \emph{single aligned pointer
store}, atomic on the architectures we target (x86-64, ARM64). A concurrent
reader therefore observes either the pre- or the post-invalidation value---never a
torn one---and, because the transition is monotonic, a stale read is benign: at
worst a loop on another thread takes the fast path for a few more entries after
an invalidation, exactly as if the redefinition had been ordered infinitesimally
later. The guard's contract is thus \emph{entry-level} consistency (the next
entry after an invalidation takes the fallback), not \emph{activation-level}
atomicity. Crucially, neither our guard nor a JIT's deoptimization makes a single
\emph{in-flight} activation atomic with respect to a concurrent redefinition;
under scalar replacement the exposure is if anything smaller than the baseline's,
because during the loop body the fast path performs only scalar arithmetic and
touches the redefinable constructor at exactly one point---the exit re-box---which
is the same \texttt{make-<t>} the unoptimized loop would call. For any
interleaving, then, the optimized loop yields a result the unoptimized loop could
also produce under some ordering of the redefinition: the transformation
introduces no new race, but neither does it manufacture a guarantee the source
program lacks. We assume redefinition is a coarse, infrequent event (a REPL or
hot-reload action) rather than high-frequency concurrent mutation, and we
evaluate only the single-threaded case; a concurrent stress evaluation is future
work.

\textbf{Novelty, and why an open world.} We claim no new escape analysis:
holding a non-escaping object's fields in variables is classical
\cite{choi1999escape,stadler2014partial}. What is new is the \emph{assembly} and
its target---loop-carried unboxing of persistent records, reached
interprocedurally by a one-level syntactic inliner and kept sound under
\emph{live class redefinition} by a per-region validity guard---in a transpiled,
latently-typed language with no optimizing IR, no JIT, and no type system to lean
on. The open world is not incidental here; it is what creates the problem the
guard solves. A JVM language such as Clojure has the identical allocation pattern
(a \texttt{defrecord} rebuilt through \texttt{loop}/\texttt{recur}) but compiles
records to classes that are not redefined mid-run, and its JIT already performs
escape analysis and deoptimizes whole frames when assumptions break; there,
neither our inliner nor our guard would earn their keep. A statically typed
functional language such as Koka attacks the same cost from the opposite
direction, with reference counting and borrowing
\cite{reinking2021perceus,lorenzen2023functional} that presuppose the very type
information FOL lacks. The gap we address---allocation-eliminating scalar
replacement for records in an \emph{open-world, untyped, JIT-less}
setting---is precisely the one those systems leave open, and it is the mechanisms
(region invalidation, safe-by-abort emit-time rewriting), not FOL itself, that
transfer to other transpiled dynamic languages in that position.

\textbf{Threats to validity.} Two are worth stating plainly. First, external
validity: FOL is a research vehicle without an outside user base, and our three
benchmarks are authored microbenchmarks that instantiate the target pattern
rather than programs sampled from a corpus (\S\ref{sec:eval}). We argue the
mechanisms transfer but do not yet demonstrate it in a second language, nor can
we yet report how often the pattern occurs in the wild. Second, internal
validity: speedups are single-machine wall-clock measurements whose baseline is
GC-scheduling sensitive; we mitigate this by repeating trials and reporting
variance alongside the deterministic allocation and native-ceiling figures, but
not by pinning the OS scheduler or sweeping microarchitectures. A corpus study
over real FOL programs and a port of the transformation to a second
persistent-first language are the two experiments that would most strengthen the
claims---and, being an emit-time syntactic rewrite rather than an IR pass, the
transformation is small enough to port.

\section{Related Work}
\label{sec:related}

\textbf{Scalar replacement and escape analysis.} Scalar replacement of
non-escaping objects is classic \cite{choi1999escape,blanchet2003escape}; the
state of the art is partial escape analysis in a JIT, which performs scalar
replacement along paths where an object does not escape and materializes it on
paths where it does \cite{stadler2014partial}. Our work adapts the core
idea---hold an object's fields in variables---to a persistent-first,
open-world dynamic language, specializing to loop accumulators and adding
run-time-redefinition soundness.

\textbf{Object inlining and aggregate scalarization.} Two further ways compilers
unbox aggregates border our work. \emph{Object inlining} fuses a referenced
object into its parent to remove an indirection and a separate allocation
\cite{dolby2000automatic}; we instead dissolve the record into loop-local scalars
rather than embed it in another heap object. \emph{Scalar replacement of
aggregates} as practiced in production compilers---LLVM's SROA, GCC's SRA---splits
a memory-resident aggregate into independent scalars once it is known not to
escape through memory. Those operate on a typed SSA form over a fixed struct
layout; ours runs on untyped surface syntax over a record layout that may change
under the running program, which is exactly what the world guard exists to
tolerate.

\textbf{Speculation, assumptions, and deoptimization.} The world guard is, in
essence, a static analogue of the assumption-and-invalidation machinery that
adaptive JIT compilers use for speculative optimization. A JIT compiles a method
under an assumption---for instance, that a class is never subclassed, licensed by
class hierarchy analysis \cite{dean1995cha}---installs a dependency on it, and
when a later class load or redefinition violates the assumption,
\emph{deoptimizes}: the optimized code is abandoned and in-flight activations are
rewritten back to interpreter frames by dynamic deoptimization
\cite{holzle1992deopt}. Production VMs track exactly these class-based
dependencies to invalidate devirtualized or inlined code when the hierarchy
changes \cite{paleczny2001hotspot,detlefs1999inlining}, and modern dynamic
compilers make speculation and its deoptimization points first-class in the IR
\cite{duboscq2013speculative}. Aggregate scalar replacement lives in the same
design space---optimize under a class-shape assumption, invalidate on
redefinition---but resolves it differently, and the differences are the
contribution. First, it is \emph{static}: instead of compiling one speculative
version and reconstructing frames on failure, it emits the fast and the fallback
path ahead of time and selects between them with a one-word guard, so there is no
interpreter to fall back to and no on-stack-replacement machinery. Second,
invalidation is \emph{region-scoped} rather than method-scoped: flipping a
class's validity cell steers only the next entry to a dependent loop, and
\emph{in-flight activations run to completion on the path they entered}---sound
precisely because the fast path performs only scalar arithmetic against the field
layout it was compiled for, so a mid-flight redefinition cannot corrupt it. The
JIT techniques buy generality (any assumption, switchable mid-execution) at the
price of a runtime, a deoptimization framework, and frame metadata; we buy a
zero-dependency emit-time rewrite at the price of switching paths only at loop
entry.

\textbf{In-place update of functional data.} Perceus achieves garbage-free
reuse via precise reference counting \cite{reinking2021perceus}; borrowing and
uniqueness extend it to ``functional but in-place'' programming
\cite{lorenzen2023functional}. Linear types \cite{wadler1990linear} give a type
discipline for single-use values. Clojure's transients \cite{hickey2009clojure}
expose in-place update as a manual, unchecked escape hatch. We target the same
allocation cost but require neither types nor programmer annotations, and we
eliminate the object rather than mutate it.

\textbf{Persistent data structures.} Hash array mapped tries
\cite{bagwell2001ideal} and their optimizations \cite{steindorfer2015optimizing}
underlie FOL's persistent collections. Deforestation \cite{wadler1990deforestation}
similarly removes intermediate structures, but by fusing producers and
consumers rather than by unboxing a loop-carried aggregate.

\section{Conclusion}

Persistent records make programs clear and allocation-heavy. We showed that a
record rebuilt every iteration of a loop can be unboxed into scalar loop
variables---no allocation on the back-edge, one re-box at the exit---by a single
classify-and-rewrite walk that reaches across a call boundary through a small
syntactic inliner and stays sound under class redefinition through a world
guard. The pass is safe by construction: it optimizes the shapes it recognizes
and declines the rest. Across four allocation-bound benchmarks it eliminates
per-iteration allocation entirely for 5.6$\times$--7.2$\times$ speedups with
identical results, recovering the allocation component of a 90--290$\times$ gap
to native mutable structs; the remainder, which we measure and attribute to
generic scalar arithmetic, is a complementary axis of specialization. We see the
transformation as the record-shaped complement to automatic transient
conversion, and a step toward bringing the benefits of partial escape analysis
to persistent-first languages with an open world.

\bibliographystyle{ACM-Reference-Format}
\bibliography{cgo2027}

\end{document}
